\begin{abstract}
LLM agents resend a growing prompt on every turn, so serving stacks reuse
cached model state for nearly every request and assume the reuse does not
change the output. We test that assumption in a pre-registered study on
llama.cpp, comparing restored state with recomputation on \TurnsPerModel{}
synthetic agent turns with tool calls, for a hybrid recurrent, a sliding-window
and a dense model. Restoring at a natural turn boundary was never bit-exact
(\NatInexact{} of \NatTurns{} turns). It changed the agent's action on
\ASparkNatShortAmbPct--\AOrnithNatLongAmbPct\% of ambiguous requests on all
three models and on none of \ExplicitRestores{} explicit-request restores;
changes concentrated where the model's top two choices were nearly tied. An
exploratory experiment points to micro-batch tiling as the cause. Aligning the
reuse boundary to the micro-batch size made restores bit-exact on the hybrid
and dense models, at a median cost of \SnapCostMin--\SnapCostMax{}
re-evaluated tokens, but not on the sliding-window model at 4{,}096-token
boundaries. Results are for synthetic traffic, one engine build, one 4-bit
quantisation, one GPU and greedy decoding.
\end{abstract}
